% ============================================================
%  sections/04_setup.tex
% ============================================================
\section{Experimental Setup}
\label{sec:setup}

\subsection{Hardware and Software}

All experiments are executed on a single CPU machine.
No GPU is used, reflecting the target deployment environment.
Software versions: Python~3.14, PyTorch~2.x, scikit-learn~1.x,
NumPy~1.x, \texttt{thop}~0.1.x.

\subsection{Dataset}

We use \textbf{Fashion-MNIST}~\citep{xiao2017fashion}, a benchmark dataset
of 28$\times$28 grayscale images of clothing articles across 10 classes
(T-shirt, Trouser, Pullover, Dress, Coat, Sandal, Shirt, Sneaker, Bag,
Ankle boot).
The dataset consists of 60{,}000 training images and 10{,}000 test images.
We hold out 10\% of training data as a validation set for
learning rate scheduling.
Pixel values are normalized to $[0,1]$.

Fashion-MNIST is preferred over MNIST because it is more
representative of real-world image classification difficulty while
remaining computationally tractable for CPU-only experiments.

\subsection{Training Hyperparameters (Baseline)}

\begin{table}[h]
\centering
\caption{Baseline training hyperparameters.}
\label{tab:hyperparams}
\small
\begin{tabular}{ll}
\toprule
Hyperparameter & Value \\
\midrule
Optimizer      & Adam~\citep{kingma2014adam} \\
Learning rate  & $10^{-3}$ \\
Weight decay   & $10^{-4}$ \\
Batch size     & 64 \\
Epochs         & 20 \\
LR scheduler   & Cosine annealing \\
Dropout        & 0.25 \\
\bottomrule
\end{tabular}
\end{table}

\subsection{Compression Hyperparameters}

\paragraph{Pruning.}
After applying the sparsity mask, we fine-tune the sparse model for
$T=5$ additional epochs with learning rate $10^{-4}$
($0.1\times$ the baseline LR) and no LR scheduler.
We evaluate at pruning ratios $\rho \in \{0.25, 0.50, 0.75, 0.90\}$.

\paragraph{Quantization.}
Dynamic INT8 PTQ requires no additional training.
We also evaluate FP32 baseline (same model, no quantization) as a
reference point.

\paragraph{Low-rank.}
We evaluate rank ratios $\rho_r \in \{0.25, 0.50, 0.75\}$.
$\rho_r = 1.0$ corresponds to the unmodified baseline.
No fine-tuning is applied after SVD approximation, as the
reconstruction is post-hoc and non-iterative.

\subsection{Evaluation Protocol}

All reported results are mean $\pm$ standard deviation across
three seeds $\{42, 123, 7\}$.
Accuracy and F1 are computed on the fixed test set (never used for
any tuning decision).
Latency statistics are computed over 100 inference runs per model
per seed, with 10 warmup runs discarded.
